% Answers to Chapter 9's exercises, lectures/L09/appendix/b_exercises.md. The book's own answers:
% the lecture publishes none. Exercises 9.3 to 9.6 are Code exercises; only the design questions
% they ask are answered, without code.
\answerchapter{9}{Structs as Drivers}

\answer{c9:ex:1}
A \code{#define}d \code{LED_ON} bakes one specific pin, and its port, into the macro: it can only
ever mean that one LED, so which LED it acts on is decided when the macro is written, not when it
is used. A \code{gpio_t} turns the pin, and the registers that go with it, into data: one set of
functions, \code{gpio_write()}, \code{gpio_toggle()} and the rest, works on whichever instance it
is handed, so which LED becomes an argument.

To support a second LED:
\begin{itemize}
  \item \textbf{Macros:} a whole new family, \code{LED2_ON}, \code{LED2_OFF}, a toggle, and its
    own setup lines, copied from the first and renamed by hand, with every copy a chance to leave a
    \code{1} where a \code{2} belongs.
  \item \textbf{\code{gpio_t}:} one more instance,
    \code{gpio_t led2 = {&DDRB, &PORTB, &PINB, 2U};}, and one more \code{gpio_init()} call (or one
    more line in \secref{c9:a:4}'s \code{leds[]} array). No new functions.
\end{itemize}

\answer{c9:ex:2}
\textbf{\code{gpio_read()} by value} would still return the right answer: the copy's \code{pinx}
holds the address of the same \code{PINx} register, so reading through it reads the same pin. What
changes is the cost and the safety net. Every call copies the whole struct instead of one 2-byte
pointer: 7 bytes for \secref{c9:a:2}'s \code{gpio_t} (three 2-byte pointers and a byte), 10 for
the demo's in \secref{c9:a:5}, on a chip with 2\,KB of RAM. And the \code{NULL} check has to go,
because a struct passed by value cannot be \code{NULL}.

\textbf{\code{gpio_init()} by value}: \secref{c9:a:2}'s version would still configure the pin
correctly, but only by accident. It never writes to the struct's own fields, only through its
pointers to the hardware registers, and the copy's pointers point at the same registers. The moment
an init function stores something in the instance, as Exercise~\ref{c9:ex:4}'s does when it
computes \code{ddrx}, \code{portx}, \code{pinx} and the bit position, the stores land in the copy,
which vanishes when the function returns, and the caller's instance is left uninitialized with no
diagnostic.

\answer{c9:ex:3}
\textbf{Pointers kept.} \code{gpio_free()} has cleared the pin's \code{DDRx} bit, so the direction
check in \code{gpio_write()} and \code{gpio_toggle()} fails and both return without touching
anything: a harmless no-op, as long as the pin stays free. The danger comes later. Once something
else claims the same pin and makes it an output again, the stale instance passes the direction
check and silently drives a pin it no longer owns.

\textbf{Pointers set to \code{NULL}.} \secref{c9:a:2}'s functions check \code{self} but not its
fields, so \code{*self->ddrx} dereferences a null pointer, which is undefined behaviour. The
ATmega328P has no memory protection to catch it: address 0 of its data space is the CPU register
\code{r0}, so instead of a crash you get a register silently read, and possibly written.

\textbf{Preferable:} set them to \code{NULL} \emph{and} add the matching check
(\code{NULL == self->ddrx}) to every function, so that a freed instance is recognisably dead and
can never reach the pin again, whoever owns it by then. Nulling the pointers without adding the
checks is worse than leaving them alone.

\answer{c9:ex:4}
No. A \code{const} member can only be given its value by an initializer where the object is
defined; it can never be assigned afterwards. In \secref{c9:a:2} that is exactly how \code{pin} got
its value, from the struct literal, so \code{const} cost nothing. Now \code{gpio_init()} computes
the bit position and assigns it (for port B, the pin number minus 8) to an object that already
exists, which a \code{const} member forbids: the compiler rejects the assignment (``assignment of
read-only member''), and casting the \code{const} away to force it is undefined behaviour. So
\code{pin} drops its \code{const}, like \code{ddrx}, \code{portx} and \code{pinx}, which never had
one. \Secref{c10:a:2} reaches the same conclusion for \code{pin} and \code{pcicrx}.

\answer{c9:ex:5}
\textbf{Why \code{static}.} At file scope \code{static} is about linkage, not lifetime: a
file-scope variable lives for the whole program either way. \code{static} gives \code{pin_reg}
internal linkage, so the name is visible only inside \file{gpio.c}, and only \code{gpio_init()} and
\code{gpio_free()} can change it. That is the whole point of the table: it enforces an invariant,
one owner per pin, and it can only do that if nothing else can edit it.

\textbf{Without \code{static}}, \code{pin_reg} has external linkage, and two things break:
\begin{itemize}
  \item Any other file can write \code{extern uint16_t pin_reg;} and set or clear bits behind the
    driver's back, handing out the same pin twice or leaking a reservation.
  \item The name is exported to the linker, where it collides with any other global
    \code{pin_reg} in the program, say a timer driver's reservation table of the same name: either
    a multiple-definition error at link time or, where uninitialized globals are merged as common
    symbols (\code{-fcommon}, GCC's default before version~10), one variable silently shared by two
    unrelated drivers.
\end{itemize}

\answer{c9:ex:6}
Yes. After \code{timer_clear()} the instance no longer owns a hardware timer/counter, and another
\code{timer_t} may already have claimed the one it had; any field left pointing at it would let the
old instance reach hardware it no longer owns, the same danger as a freed \code{gpio_t} in
Exercise~\ref{c9:ex:3}. The safe, inert values are the ones every function can recognise as ``owns
nothing'':
\begin{itemize}
  \item The reserved timer/counter set to a value no timer/counter has (e.g.\ \code{0xFF}), with
    every function checking it first, so \code{timer_start()} or \code{timer_toggle()} on a cleared
    instance does nothing instead of restarting hardware someone else now owns.
  \item Running \code{false} and the elapsed-time count \code{0}, so \code{timer_running()} reports
    it as stopped.
  \item The callback \code{NULL}, so nothing can call through it, and the timeout \code{0}, the one
    value \code{timer_init()} refuses, so the instance reads as never initialized.
\end{itemize}
The hardware side goes first: stop the timer/counter's clock, disable its compare-match interrupt,
and clear the driver's own record of which instance owns it, before the reservation bit is
released, so that its interrupt can never fire for an instance that no longer owns it.
